% GENERATED FILE. Edit canonical lessons, then run npm run generate:books.
% canonical-source-hash: fnv1a64:1e1bd80ca7fc0636
% canonical-lessons: UR-C459-fayyaz, UR-C459-khud-gharaz, UR-C459-sakhi, UR-C459-bakhil, UR-C459-sadqa

\chapter{Generous and Mean}
\label{ch:ur-generous-and-mean}
\hlchaptermodality{459}

\begin{chapteropening}
\textbf{By the end of this chapter:} \emph{I can say whether someone is generous or mean.}

Complete the last lesson of chapter 459: I can say whether someone is generous or mean.
\end{chapteropening}

\section[\texorpdfstring{\ur{فیاض} (fayyāz)}{فیاض (fayyāz)}]{\ur{فیاض} (fayyāz) — generous; bountiful}

\label{lesson:UR-C459-fayyaz}



\begin{quote}
Before the new one: say the Urdu for to snap at, then the Urdu for to bear.
\end{quote}

\subsection*{You'll want to know: \ur{فیاض}}
\textbf{\ur{فیاض}} — \emph{fayyāz} — \textquotedblleft{}generous; bountiful\textquotedblright{}.

\begin{quote}
\ur{ہمارے} \ur{پڑوسی} \ur{بڑے} \ur{فیاض} \ur{ہیں۔}
\end{quote}

\begin{quote}
\emph{hamāre paṛosī baṛe fayyāz haiṅ.}
\end{quote}

\begin{quote}
Our neighbours are very generous.
\end{quote}

\begin{cousinweb}
\textbf{\ur{فیاض}} (\emph{fayyāz}) is usually traced, through Persian, to Arabic \emph{fayyāḍ}, \textquotedblleft{}overflowing, bountiful\textquotedblright{}.
\end{cousinweb}

\begin{grammarlens}[title={What you've built}]
One more word for describing what you lived and felt.
\end{grammarlens}

\subsection*{Your turn}
\begin{itemize}
\raggedright
  \item \emph{Say it:} \emph{fayyāz}
  \item \emph{Say it:} \emph{fayyāz}, once more
  \item \emph{Say it:} \emph{hamāre paṛosī baṛe fayyāz haiṅ.} — Our neighbours are very generous.
  \item \emph{Recall:} say the Urdu for to snap at, then the Urdu for to bear, then say \emph{fayyāz} again
\end{itemize}

\begin{tcolorbox}[breakable,colback=teal!4,colframe=teal!35!black,title={Before you move on}]
What does \ur{فیاض} mean? (\textquotedblleft{}Generous; bountiful\textquotedblright{}.) Say it once more.
\end{tcolorbox}
\section[\texorpdfstring{\ur{خود} \ur{غرض} (khud gharaz)}{خود غرض (khud gharaz)}]{\ur{خود} \ur{غرض} (khud gharaz) — selfish}

\label{lesson:UR-C459-khud-gharaz}



\begin{quote}
Before the new one: say the Urdu for to bear, then the Urdu for bountiful.
\end{quote}

\subsection*{You'll want to know: \ur{خود} \ur{غرض}}
\textbf{\ur{خود} \ur{غرض}} — \emph{khud gharaz} — \textquotedblleft{}selfish\textquotedblright{}.

\begin{quote}
\ur{وہ} \ur{خود} \ur{غرض} \ur{ہے،} \ur{کسی} \ur{کی} \ur{مدد} \ur{نہیں} \ur{کرتا۔}
\end{quote}

\begin{quote}
\emph{voh khud gharaz hai, kisī kī madad nahīṅ kartā.}
\end{quote}

\begin{quote}
He is selfish and helps nobody.
\end{quote}

\begin{cousinweb}
\textbf{\ur{خود} \ur{غرض}} (\emph{khud gharaz}) is built from \textbf{\ur{خود}}, \textquotedblleft{}oneself\textquotedblright{}, usually traced to Persian \emph{khud}, and \emph{gharaz}, \textquotedblleft{}an aim\textquotedblright{}, usually traced to Arabic \emph{gharaḍ}: someone whose aim is himself.
\end{cousinweb}

\begin{grammarlens}[title={What you've built}]
One more word for describing what you lived and felt.
\end{grammarlens}

\subsection*{Your turn}
\begin{itemize}
\raggedright
  \item \emph{Say it:} \emph{khud gharaz}
  \item \emph{Say it:} \emph{khud gharaz}, once more
  \item \emph{Say it:} \emph{voh khud gharaz hai, kisī kī madad nahīṅ kartā.} — He is selfish and helps nobody.
  \item \emph{Recall:} say the Urdu for to bear, then the Urdu for bountiful, then say \emph{khud gharaz} again
\end{itemize}

\begin{tcolorbox}[breakable,colback=teal!4,colframe=teal!35!black,title={Before you move on}]
What does \ur{خود} \ur{غرض} mean? (\textquotedblleft{}Selfish\textquotedblright{}.) Say it once more.
\end{tcolorbox}
\section[\texorpdfstring{\ur{سخی} (sakhī)}{سخی (sakhī)}]{\ur{سخی} (sakhī) — open-handed}

\label{lesson:UR-C459-sakhi}



\begin{quote}
Before the new one: say the Urdu for bountiful, then the Urdu for selfish.
\end{quote}

\subsection*{You'll want to know: \ur{سخی}}
\textbf{\ur{سخی}} — \emph{sakhī} — \textquotedblleft{}open-handed\textquotedblright{}.

\begin{quote}
\ur{سخی} \ur{آدمی} \ur{سب} \ur{کی} \ur{مدد} \ur{کرتا} \ur{ہے۔}
\end{quote}

\begin{quote}
\emph{sakhī ādmī sab kī madad kartā hai.}
\end{quote}

\begin{quote}
An open-handed man helps everyone.
\end{quote}

\begin{cousinweb}
\textbf{\ur{سخی}} (\emph{sakhī}) is usually traced, through Persian, to Arabic \emph{sakhī}, \textquotedblleft{}generous\textquotedblright{}.
\end{cousinweb}

\begin{grammarlens}[title={What you've built}]
One more word for describing what you lived and felt.
\end{grammarlens}

\subsection*{Your turn}
\begin{itemize}
\raggedright
  \item \emph{Say it:} \emph{sakhī}
  \item \emph{Say it:} \emph{sakhī}, once more
  \item \emph{Say it:} \emph{sakhī ādmī sab kī madad kartā hai.} — An open-handed man helps everyone.
  \item \emph{Recall:} say the Urdu for bountiful, then the Urdu for selfish, then say \emph{sakhī} again
\end{itemize}

\begin{tcolorbox}[breakable,colback=teal!4,colframe=teal!35!black,title={Before you move on}]
What does \ur{سخی} mean? (\textquotedblleft{}Open-handed\textquotedblright{}.) Say it once more.
\end{tcolorbox}
\section[\texorpdfstring{\ur{بخیل} (bakhīl)}{بخیل (bakhīl)}]{\ur{بخیل} (bakhīl) — miserly}

\label{lesson:UR-C459-bakhil}



\begin{quote}
Before the new one: say the Urdu for selfish, then the Urdu for open-handed.
\end{quote}

\subsection*{You'll want to know: \ur{بخیل}}
\textbf{\ur{بخیل}} — \emph{bakhīl} — \textquotedblleft{}miserly\textquotedblright{}.

\begin{quote}
\ur{وہ} \ur{بخیل} \ur{آدمی} \ur{پیسہ} \ur{خرچ} \ur{نہیں} \ur{کرتا۔}
\end{quote}

\begin{quote}
\emph{voh bakhīl ādmī paisā kharch nahīṅ kartā.}
\end{quote}

\begin{quote}
That miserly man does not spend money.
\end{quote}

\begin{cousinweb}
\textbf{\ur{بخیل}} (\emph{bakhīl}) is usually traced, through Persian, to Arabic \emph{bakhīl}, \textquotedblleft{}miserly\textquotedblright{}.
\end{cousinweb}

\begin{grammarlens}[title={What you've built}]
One more word for describing what you lived and felt.
\end{grammarlens}

\subsection*{Your turn}
\begin{itemize}
\raggedright
  \item \emph{Say it:} \emph{bakhīl}
  \item \emph{Say it:} \emph{bakhīl}, once more
  \item \emph{Say it:} \emph{voh bakhīl ādmī paisā kharch nahīṅ kartā.} — That miserly man does not spend money.
  \item \emph{Recall:} say the Urdu for selfish, then the Urdu for open-handed, then say \emph{bakhīl} again
\end{itemize}

\begin{tcolorbox}[breakable,colback=teal!4,colframe=teal!35!black,title={Before you move on}]
What does \ur{بخیل} mean? (\textquotedblleft{}Miserly\textquotedblright{}.) Say it once more.
\end{tcolorbox}
\section[\texorpdfstring{\ur{صدقہ} (sadqā)}{صدقہ (sadqā)}]{\ur{صدقہ} (sadqā) — alms (given for a blessing)}

\label{lesson:UR-C459-sadqa}



\begin{quote}
Before the new one: say the Urdu for open-handed, then the Urdu for miserly.
\end{quote}

\subsection*{You'll want to know: \ur{صدقہ}}
\textbf{\ur{صدقہ}} — \emph{sadqā} — \textquotedblleft{}alms (given for a blessing)\textquotedblright{}.

\begin{quote}
\ur{دادی} \ur{ہر} \ur{ہفتے} \ur{صدقہ} \ur{دیتی} \ur{ہیں۔}
\end{quote}

\begin{quote}
\emph{dādī har hafte sadqā detī haiṅ.}
\end{quote}

\begin{quote}
My grandmother gives alms every week.
\end{quote}

Masculine.

\begin{cousinweb}
\textbf{\ur{صدقہ}} (\emph{sadqā}) is usually traced, through Persian, to Arabic \emph{ṣadaqa}, \textquotedblleft{}alms\textquotedblright{}.
\end{cousinweb}

\begin{grammarlens}[title={What you've built}]
One more word for describing what you lived and felt.
\end{grammarlens}

\subsection*{Your turn}
\begin{itemize}
\raggedright
  \item \emph{Say it:} \emph{sadqā}
  \item \emph{Say it:} \emph{sadqā}, once more
  \item \emph{Say it:} \emph{dādī har hafte sadqā detī haiṅ.} — My grandmother gives alms every week.
  \item \emph{Recall:} say the Urdu for open-handed, then the Urdu for miserly, then say \emph{sadqā} again
\end{itemize}

\begin{tcolorbox}[breakable,colback=teal!4,colframe=teal!35!black,title={Before you move on}]
What does \ur{صدقہ} mean? (\textquotedblleft{}Alms (given for a blessing)\textquotedblright{}.) Say it once more.
\end{tcolorbox}
